% PAGES: 60-61  (printed 44-45)
% Continues section 2.3 "Backward Substitution", begun before page 44
% (PDF page 60). Do not open or append to the file containing the start
% of section 2.3 -- start reading from here.

By multiplying the $n$-th row of $U$ by $x$, and using the fact that
$U(n,k) = 0$ for all $k = 1:(n-1)$, since $U$ is upper triangular, see
(2.13), we obtain that
\[
\sum_{k=1}^{n} U(n,k)x(k) \;=\; U(n,n)x(n) \;=\; b(n),
\]
and therefore
\begin{equation}
x(n) \;=\; \frac{b(n)}{U(n,n)};
\end{equation}
note that $U(n,n) \neq 0$, since $U$ is nonsingular.

Moving backward to computing $x(n-1)$, multiply the $(n-1)$-th row of $U$
by $x$, and use the fact that $U(n-1,k) = 0$ for all $k = 1:(n-2)$, since
$U$ is upper triangular, see (2.13), to obtain that
\begin{equation}
\sum_{k=1}^{n} U(n-1,k)x(k) \;=\; U(n-1,n-1)x(n-1) + U(n-1,n)x(n) \;=\; b(n-1).
\end{equation}

By solving (2.15) for $x(n-1)$, we find that
\begin{equation}
x(n-1) \;=\; \frac{b(n-1) - U(n-1,n)x(n)}{U(n-1,n-1)};
\end{equation}
note that $U(n-1,n-1) \neq 0$, since $U$ is nonsingular. Since the value
of $x(n)$ is known from (2.14), the value of $x(n-1)$ can be computed
from (2.16).

We continue to move backward until all the entries of $x$ are computed,
as follows: Assume that the values of $x(n), x(n-1), \ldots, x(j+1)$
have been computed, where $n-1 \geq j \geq 1$. To compute $x(j)$,
multiply the $j$-th row of $U$ by $x$, and use the fact that $U(j,k) = 0$
for all $k = 1:(j-1)$, since $U$ is upper triangular, to obtain that
\[
\sum_{k=1}^{n} U(j,k)x(k) \;=\; \sum_{k=j}^{n} U(j,k)x(k) \;=\; b(j).
\]
Thus,
\begin{equation}
U(j,j)x(j) \;+\; \sum_{k=j+1}^{n} U(j,k)x(k) \;=\; b(j).
\end{equation}
By solving (2.17) for $x(j)$, we find that
\begin{equation}
x(j) \;=\; \frac{b(j) - \sum_{k=j+1}^{n} U(j,k)x(k)}{U(j,j)};
\end{equation}
note that $U(j,j) \neq 0$, since $L$ is nonsingular. Thus, given
$x(n), x(n-1), \ldots, x(j+1)$, the value for $x(j)$ can be computed
using (2.18).

The pseudocode for Backward Substitution can be found in Table 2.3.

The operation count for Backward Substitution is $n^2 + O(n)$, the same
as for Forward Substitution; the proof of the result below is similar to
that of Lemma 2.1:

\begin{lemma}
The operation count for the Backward Substitution algorithm applied to
an upper triangular matrix of size $n$ is
\begin{equation}
n^2 \;+\; O(n).
\end{equation}
\end{lemma}

\begin{table}[H]
\centering
\caption{Pseudocode for Backward Substitution}
\fbox{%
\begin{minipage}{0.82\textwidth}
Function Call:\\
$x = \text{backward\_subst}(U,b)$

\medskip
Input:\\
$U = $ nonsingular upper triangular matrix of size $n$\\
$b = $ column vector of size $n$

\medskip
Output:\\
$x = $ solution to $Ux = b$

\medskip
$x(n) = \dfrac{b(n)}{U(n,n)}$;\\
for $j = (n-1):1$\\
\hspace*{1.5em}sum $= 0$;\\
\hspace*{1.5em}for $k = (j+1):n$\\
\hspace*{3em}sum $=$ sum $+ U(j,k)x(k)$;\\
\hspace*{1.5em}end\\
\hspace*{1.5em}$x(j) = \dfrac{b(j)-\text{sum}}{U(j,j)}$;\\
end
\end{minipage}%
}
\end{table}

The relevance of Forward Substitution and Backward Substitution is
highlighted by the fact that, for a nonsingular matrix $A$ that is not
upper triangular nor lower triangular, a solution to the linear system
$Ax = b$ can be found without explicitly computing $A^{-1}$, by using
forward substitution and backward substitution, and the LU decomposition
with row pivoting of $A$; see section 2.6 and the pseudocode from
Table 2.10 for details.

We conclude this section with the Backward Substitution pseudocode
corresponding to an upper triangular bidiagonal matrix, which will be
used for solving linear systems corresponding to tridiagonal matrices;
cf. section 2.5.1 and section 6.3.

Let $U$ be an $n \times n$ upper triangular bidiagonal matrix, i.e., a
matrix whose only nonzero entries could be the main diagonal entries
$U(j,j)$, for $j = 1:n$, and the upper diagonal entries $U(j,j+1)$, for
$j = 1:(n-1)$. Then, the recursion formula (2.18) simplifies to
\[
x(j) \;=\; \frac{b(j) - U(j,j+1)x(j+1)}{U(j,j)},
\]
which corresponds to the reduced version of the Backward Substitution
pseudocode from Table 2.3 included in Table 2.4.

The operation count for the Backward Substitution for upper triangular
bidiagonal matrices is
\begin{equation}
1 + 3(n-1) \;=\; 3n - 2.
\end{equation}
